\section{The Binding: One Commitment for Four Proofs}
\label{sec:binding}

Before specifying \emph{how} a worker proves it ran a model
(Sections~\ref{sec:four-proofs}--\ref{sec:freivalds}), we fix \emph{what} the
proof is about. The binding is a single $32$-byte value, $\mathsf{reportData}$,
that names the task, the measured model, the input, the output, and the runtime
under which the computation must have occurred. Every proof system---TEE,
optimistic-Freivalds, zkML, or $M$-of-$N$---attests to the same
$\mathsf{reportData}$; the proof system is a backend, and $\mathsf{reportData}$ is
the join key. This is the decomplecting move: \emph{what was computed} (the
binding) is separated from \emph{how we believe it} (the backend), so the two vary
independently.

\subsection{Domain-separated two-stage construction}

The binding is built in two keccak stages with disjoint domain tags, mirroring the
deployed \texttt{ComputeProofLib} library. Both tags are raw UTF-8 byte strings
used as keccak prefixes (no length framing), and both are versioned:
\begin{align}
\mathsf{DOMAIN\_CHALLENGE} &= \texttt{"lux/aivm/compute-challenge/v1"}, \\
\mathsf{DOMAIN\_REPORT}    &= \texttt{"lux/aivm/compute-report/v1"}.
\end{align}
The \emph{challenge} commits to everything fixed \emph{before} the worker runs:
\begin{equation}
\label{eq:challenge}
\mathsf{ch} \;=\; \mathsf{keccak}\Bigl(
  \mathsf{DOMAIN\_CHALLENGE} \,\|\,
  \mathsf{taskId} \,\|\,
  \mathsf{intentID} \,\|\,
  \mathsf{modelSpecHash} \,\|\,
  \mathsf{inputHash} \,\|\,
  \mathsf{openBlockHash} \,\|\,
  \mathsf{operator}
\Bigr),
\end{equation}
and the \emph{report} commits the challenge together with the produced output and
the runtime under which it was produced:
\begin{equation}
\label{eq:report}
\mathsf{reportData} \;=\; \mathsf{keccak}\Bigl(
  \mathsf{DOMAIN\_REPORT} \,\|\,
  \mathsf{ch} \,\|\,
  \mathsf{workloadType} \,\|\,
  \mathsf{modelSpecHash} \,\|\,
  \mathsf{inputHash} \,\|\,
  \mathsf{outputHash} \,\|\,
  \mathsf{runtimeMeasurement}
\Bigr).
\end{equation}
The wire encoding is byte-pinned so that the Solidity precompile, the Go
attestation precompile, and the Rust engine produce identical digests:
$\mathsf{taskId}$ is a $32$-byte big-endian integer; $\mathsf{intentID}$,
$\mathsf{modelSpecHash}$, $\mathsf{inputHash}$, $\mathsf{openBlockHash}$,
$\mathsf{outputHash}$, and $\mathsf{runtimeMeasurement}$ are $32$-byte words; and
$\mathsf{operator}$ is a $20$-byte address, \emph{not} left-padded. The deployed
library realizes~\eqref{eq:challenge}--\eqref{eq:report} verbatim over
\texttt{abi.encodePacked}; this paper adds the single field $\mathsf{workloadType}$
to~\eqref{eq:report} (Section~\ref{sec:four-proofs}) so that one binding serves all
four proofs without overloading any other field.

\subsection{The fields, and what each defends}

Each field in~\eqref{eq:challenge}--\eqref{eq:report} neutralizes a specific cheat.

\begin{description}[leftmargin=0pt,style=nextline]
\item[$\mathsf{modelSpecHash}$ --- a \emph{measured} weights Merkle root, not a
label.] This is the keccak Merkle root over the loaded, quantized model bytes
(Section~\ref{sec:freivalds}), \emph{computed at load time from the weights the
worker actually holds in memory}, not a human-readable model name. Binding the
measured root defeats \emph{cheap-model substitution}: a worker cannot answer with
a $1$B model while claiming the fee schedule of a $70$B model, because the proof
backend re-derives activations against the committed root and a different model
yields a different root and different activations.
\item[$\mathsf{inputHash}$.] Commits the exact prompt/embedding input. A worker
cannot answer an easier question than the one asked.
\item[$\mathsf{outputHash}$.] Commits the produced output (token ids for
inference, the embedding vector for embedding, the LoRA-delta root for training).
This is what the proof must show is the genuine result.
\item[$\mathsf{openBlockHash}$ --- the unpredictability beacon.] The hash of a
recent block that is \emph{not known when the task is posted}. Because $\mathsf{ch}$
mixes in $\mathsf{openBlockHash}$, the challenge---and hence the random vector the
Freivalds check derives from it (Section~\ref{sec:freivalds})---cannot be predicted
or pre-baked. This is the anti-replay and anti-precomputation root.
\item[$\mathsf{operator}$ --- per-operator domain.] The provider's address.
Because each operator's $\mathsf{ch}$ and $\mathsf{reportData}$ are distinct, a
proof produced by operator $A$ does not verify for operator $B$; honest work cannot
be copied and re-sold under another identity (the same anti-copy property the
A-Chain commitment of Section~\ref{sec:guess-to-mint} sought, now extended to the
\emph{proof} rather than only the output).
\item[$\mathsf{runtimeMeasurement}$ --- the determinism pin.] A digest of the
runtime and sampler configuration under which the output is reproducible: the
kernel set, tokenizer, quantization mode, temperature, and seed
(Section~\ref{sec:freivalds}). Pinning it makes ``re-derive and compare'' a
well-defined operation across heterogeneous hardware.
\end{description}

\subsection{Security of the binding}

We model keccak as a random oracle $\mathcal{H}$ and quantify the two guarantees
the construction must provide: a forged report cannot be replayed onto a different
task, and it cannot be silently re-pointed to a different model or input. Both
reduce to collision/preimage resistance, but the domain separation and the beacon
make the reductions tight and the attack windows closed.

\begin{lemma}[Anti-substitution]
\label{lem:anti-sub}
Model $\mathcal{H}$ as a random oracle. If two reports share a $\mathsf{reportData}$
value but differ in any of $\mathsf{modelSpecHash}$, $\mathsf{inputHash}$,
$\mathsf{outputHash}$, $\mathsf{workloadType}$, or (through $\mathsf{ch}$)
$\mathsf{taskId}$, $\mathsf{intentID}$, $\mathsf{openBlockHash}$, or
$\mathsf{operator}$, then the adversary has computed a collision of $\mathcal{H}$.
After $Q$ oracle queries the probability of producing any such pair is at most
$\binom{Q}{2} 2^{-256} \le Q^2 2^{-257}$.
\end{lemma}

\begin{proof}
$\mathsf{reportData} = \mathcal{H}(x)$ where $x$ is the encoding
of~\eqref{eq:report}, which contains $\mathsf{ch} = \mathcal{H}(y)$ with $y$ the
encoding of~\eqref{eq:challenge}. The packed encoding is injective on the listed
fields: every field has a fixed byte length ($\mathsf{operator}$ is $20$ bytes, all
other components are $32$ bytes, $\mathsf{workloadType}$ is a fixed-width tag) and
the domain prefixes are constants of fixed length, so distinct field tuples yield
distinct byte strings $x$ (resp.\ $y$). Two reports with equal $\mathsf{reportData}$
and unequal $x$ are a direct collision of $\mathcal{H}$; if $x$ is equal but the
challenge inputs differ, then $\mathsf{ch}$ collided, i.e.\ a collision of
$\mathcal{H}$ on $y$. A random oracle answers each fresh query with a uniform
$256$-bit string, so any fixed pair of distinct inputs collides with probability
$2^{-256}$; a union bound over the $\binom{Q}{2}$ query pairs gives the stated
bound. \qed
\end{proof}

The disjoint domain tags do real work in Lemma~\ref{lem:anti-sub}: without them a
$\mathsf{challenge}$ digest and a $\mathsf{reportData}$ digest could in principle be
made to coincide and be cross-used (a type-confusion between the two stages).
Prefixing each stage with its own constant forces the two oracle ``slices'' apart,
so a challenge can never be mistaken for a report and vice versa. The same
construction generalizes to any future proof family by adding a new
$\mathsf{workloadType}$ tag rather than a new field, keeping the binding stable.

\begin{lemma}[Anti-replay / freshness]
\label{lem:anti-replay}
Suppose $\mathsf{openBlockHash}$ is the hash of a block whose contents are
unpredictable until height $H_{\mathrm{open}}$, and that a task posted at height
$H_{\mathrm{post}} < H_{\mathrm{open}}$ fixes $\mathsf{openBlockHash} :=
\mathsf{blockhash}(H_{\mathrm{open}})$. Then any $\mathsf{reportData}$, and any
random challenge vector derived from it, that an adversary commits to before height
$H_{\mathrm{open}}$ is independent of the value that will be required, except with
the probability that the adversary predicts $\mathsf{blockhash}(H_{\mathrm{open}})$.
Under the random-oracle model for the block-hash beacon this prediction probability
is $2^{-256}$ per guess.
\end{lemma}

\begin{proof}
$\mathsf{ch}$ is a function of $\mathsf{openBlockHash}$ through $\mathcal{H}$
(Eq.~\ref{eq:challenge}), and $\mathsf{reportData}$ is a function of $\mathsf{ch}$
through $\mathcal{H}$ (Eq.~\ref{eq:report}); the Freivalds vector is in turn derived
from $\mathsf{ch}$ (Section~\ref{sec:freivalds}). A value committed before
$H_{\mathrm{open}}$ is a function of information available before
$H_{\mathrm{open}}$, which by hypothesis is independent of
$\mathsf{blockhash}(H_{\mathrm{open}})$. By the random-oracle property of the beacon,
the only way a pre-$H_{\mathrm{open}}$ commitment can equal the post-$H_{\mathrm{open}}$
requirement is if it embeds a correct guess of the unpredictable beacon output, an
event of probability $2^{-256}$ per attempt. \qed
\end{proof}

Lemma~\ref{lem:anti-replay} is what prevents a prover from \emph{pre-baking} a
proof: it cannot precompute activations or a Freivalds response for a challenge it
does not yet know, and it cannot replay a stale proof, because the beacon ties the
challenge to a block that did not exist when a stale proof was made. Together with
Lemma~\ref{lem:anti-sub} and the per-operator field, the binding guarantees that a
valid proof is \emph{about this task, this model, this input, this operator, and
this freshly-sampled challenge}---the necessary frame within which
Section~\ref{sec:freivalds}'s soundness theorem then guarantees the computation
itself was performed.
